% Relatório técnico — Trabalho Avaliativo Agentes Inteligentes
% Formatação alinhada à ABNT NBR 14724: margens, fonte 12 pt, espaçamento 1,5, recuo 1,25 cm
% Compilar: pdflatex trabalho_agentes.tex (duas vezes para sumário)
\documentclass[12pt,a4paper,oneside]{article}

\usepackage[T1]{fontenc}
\usepackage[utf8]{inputenc}
\usepackage[brazil]{babel}
\usepackage{mathptmx}

\usepackage[a4paper,left=3cm,right=2cm,top=3cm,bottom=2cm]{geometry}

\usepackage{setspace}
\onehalfspacing

\usepackage{indentfirst}
\setlength{\parindent}{1.25cm}

\usepackage{hyperref}
\hypersetup{colorlinks=true,linkcolor=black,urlcolor=blue,citecolor=black}
% Evita âncora duplicada "page.1" (capa + reinício da numeração)
\AtBeginDocument{\hypersetup{pageanchor=false}}

\usepackage{graphicx}
\usepackage{booktabs}
\usepackage{enumitem}
\setlist{nosep,leftmargin=1.25cm}

\usepackage{titlesec}
\titleformat{\section}{\normalfont\bfseries\fontsize{12}{18}\selectfont}{\thesection}{1em}{}
\titlespacing*{\section}{0pt}{18pt}{12pt}
\titleformat{\subsection}{\normalfont\bfseries\fontsize{12}{18}\selectfont}{\thesubsection}{1em}{}
\titlespacing*{\subsection}{0pt}{12pt}{6pt}

\renewcommand{\refname}{REFERÊNCIAS}

\begin{document}

% --- Capa (ABNT: dados da instituição, título, autor, local, ano) ---
\begin{titlepage}
\centering
\onehalfspacing
\vspace*{1.5cm}

{\fontsize{14}{16}\selectfont\textbf{INSTITUTO DE ENSINO SUPERIOR}}\\[0.3cm]
{\fontsize{12}{14}\selectfont Pós-Graduação em Inteligência Artificial}\\[0.5cm]
{\fontsize{12}{14}\selectfont Disciplina: Agentes Inteligentes}\\[0.5cm]
{\fontsize{12}{14}\selectfont Período letivo: 2026}

\vfill

{\fontsize{14}{18}\selectfont\textbf{SIMULAÇÃO DE ORQUESTRA DE DRONES EM GRADE:\\INTEGRAÇÃO DE GRANDES MODELOS DE LINGUAGEM\\E AGENTE BASEADO EM MODELO}}\\[0.8cm]
{\fontsize{12}{14}\selectfont Trabalho prático --- Simulação e agência inteligente}

\vfill

\begin{flushleft}
\fontsize{12}{14}\selectfont
\textbf{Autor:}\\[0.4cm]
Robson Medeiros Santos\\
\textit{Contribuição:} concepção e implementação integral do projeto: modelagem do ambiente em grade (\texttt{environment.py}), agente coreógrafo com integração ao LLM via Groq (\texttt{llm\_client.py}, \texttt{agents/choreographer.py}), agente piloto baseado em modelo com BFS (\texttt{agents/pilot.py}), simulação e orquestração (\texttt{simulation.py}, \texttt{main.py}), interface gráfica opcional (\texttt{gui.py}), documentação no repositório e elaboração deste relatório.
\end{flushleft}

\vfill

{\fontsize{12}{14}\selectfont Teresina -- PI}\\
{\fontsize{12}{14}\selectfont 2026}

\end{titlepage}
\hypersetup{pageanchor=true}

\newpage
\setcounter{page}{1}

% --- Resumo ---
\section*{RESUMO}
\addcontentsline{toc}{section}{RESUMO}

Este trabalho apresenta uma simulação original de \textit{show} de drones em ambiente de grade discreta, desenvolvida em linguagem Python sem dependências externas obrigatórias (\textit{pip}). Dois agentes autônomos operam de forma coordenada: um \textbf{coreógrafo}, cuja decisão de alto nível é obtida por meio de um grande modelo de linguagem (LLM) via API Groq, com engenharia de \textit{prompt} e saída estruturada em JSON; e um \textbf{piloto}, implementado como agente \textbf{baseado em modelo}, que planeja deslocamentos com busca em largura (BFS), evitando obstáculos e colisões entre aeronaves. O ambiente inclui múltiplos obstáculos, resolução de alvos inviáveis e modo de demonstração quando a API não está disponível. O repositório público contém o código-fonte e instruções de execução. O documento formaliza a especificação PEAS e a caracterização do ambiente segundo as dimensões de Russell e Norvig.

\vspace{0.5cm}
\noindent\textbf{Palavras-chave:} agentes inteligentes; grandes modelos de linguagem; simulação multiagente; planejamento baseado em modelo; Python.

\newpage
\tableofcontents
\newpage

% --- Corpo textual ---
\section{INTRODUÇÃO}

A disciplina de Agentes Inteligentes exige a construção de uma simulação em que pelo menos dois agentes autônomos interajam com o meio ou entre si, sendo obrigatório que ao menos um deles utilize raciocínio baseado em LLM e outro siga uma das arquiteturas clássicas apresentadas por Russell e Norvig \cite{russell2021}.

O tema escolhido --- orquestração de drones em grade, inspirada em formações de \textit{drone show} --- oferece um contexto visual e conceitualmente rico: o \textbf{coreógrafo} interpreta o estado do espetáculo e define formações; o \textbf{piloto} executa trajetórias seguras. A escolha reduz a repetição de temas frequentes (armazém, predador-presa genérico) e mantém clara a separação entre deliberação simbólica (LLM) e planejamento geométrico (modelo interno + BFS).

\section{IDENTIFICAÇÃO DO AUTOR E CONTRIBUIÇÃO}

Conforme solicitado no enunciado, segue a identificação do autor e a descrição objetiva da contribuição no projeto.

\begin{table}[h]
\centering
\small
\begin{tabular}{p{4.2cm}p{9.5cm}}
\toprule
\textbf{Autor} & \textbf{Contribuição no projeto} \\
\midrule
Robson Medeiros Santos & Modelagem do ambiente em grade com obstáculos e formações (\texttt{environment.py}); coreógrafo com LLM via API Groq e saída estruturada (\texttt{llm\_client.py}, \texttt{agents/choreographer.py}); piloto baseado em modelo com BFS e anti-colisão (\texttt{agents/pilot.py}); laço de simulação e CLI (\texttt{simulation.py}, \texttt{main.py}); interface gráfica Tkinter (\texttt{gui.py}); repositório público, testes de execução e relatório em conformidade com a NBR 14724. \\
\bottomrule
\end{tabular}
\end{table}

\section{ESPECIFICAÇÃO DA TAREFA (PEAS)}

A análise PEAS (\textit{Performance measure}, \textit{Environment}, \textit{Actuators}, \textit{Sensors}) descreve formalmente a tarefa de cada agente em relação ao ambiente simulado \cite{russell2021}.

\subsection{Medida de desempenho (\textit{Performance})}

Para o sistema como um todo, considera-se desempenho satisfatório quando: (i) os quatro drones ocupam, em conjunto, o conjunto de células-alvo definido pela formação vigente, sem colisões entre si e sem violação dos limites da grade; (ii) obstáculos não são atravessados; (iii) o coreógrafo só redefine a formação após a \textbf{conclusão} da anterior (transição detectada quando o multiconjunto de posições coincide com o multiconjunto de alvos), preservando coerência narrativa do ``espetáculo''; (iv) opcionalmente, minimiza-se o número de passos até a convergência, sem ser critério único.

\subsection{Ambiente (\textit{Environment})}

O ambiente é uma grade retangular discreta de dimensões $10 \times 14$ células, com conjunto fixo de obstáculos (paredes e blocos) e quatro drones identificados. O tempo é discretizado em \textit{beats}: a cada passo, todos os drones aplicam simultaneamente um movimento escolhido. O estado inclui posições dos drones, parâmetros da formação atual (nome, centro, escala), obstáculos e contador de tempo. Não há fenômenos físicos contínuos: posições são inteiras.

\subsection{Atuadores (\textit{Actuators})}

Cada drone pode executar, por turno, um entre os comandos: deslocamento para norte, sul, leste ou oeste, ou permanência (\textsc{hold}). O coreógrafo não move drones diretamente: altera parâmetros globais da formação (tipo geométrico, centro aproximado, escala), que o ambiente traduz em células-alvo (com correção para obstáculos e unicidade de alvo).

\subsection{Sensores (\textit{Sensors})}

Na implementação, o coreógrafo recebe um resumo textual do estado (posições, obstáculos, \textit{beat}, formação corrente), correspondendo a observabilidade total do simulador para esse agente. O piloto utiliza o mesmo modelo de mundo disponível no código (posições, obstáculos, alvos alinhados por proximidade), compatível com a arquitetura baseada em modelo, em que o agente mantém representação explícita do espaço e atualiza planos conforme o estado.

\section{ANÁLISE DO AMBIENTE SEGUNDO RUSSELL E NORVIG}

Caracteriza-se o simulador segundo as seis dimensões usuais de análise de ambientes de tarefas \cite{russell2021}.

\subsection{Totalmente observável versus parcialmente observável}

O ambiente é \textbf{totalmente observável} para os agentes tal como implementados: todo o estado relevante (grade, obstáculos, drones, parâmetros de formação) está disponível à função de decisão, sem incerteza sensorial modelada.

\subsection{Agente único versus multiagente}

Trata-se de ambiente \textbf{multiagente}: quatro drones competem por espaço com regras de precedência no deslocamento; o coreógrafo e o piloto são dois processos decisórios distintos que influenciam o mesmo estado. A interação é cooperativa (objetivo comum de formação), com restrições de não colisão.

\subsection{Determinístico versus estocástico}

O motor de simulação e o piloto são \textbf{determinísticos} para um dado estado e decisão do coreógrafo. O núcleo do coreógrafo via LLM introduz \textbf{estocasticidade} nas escolhas de formação (temperatura e variabilidade do modelo). O modo demonstração (\textit{stub}) é determinístico condicionado ao \textit{beat}.

\subsection{Episódico versus sequencial}

O problema é \textbf{sequencial}: a decisão corrente afeta estados futuros (posições, viabilidade de caminhos, cumprimento da formação). Não há independência entre ``turnos'' no sentido episódico estrito.

\subsection{Estático versus dinâmico}

O ambiente é \textbf{dinâmico}: enquanto um drone delibera no modelo simplificado, outros movem-se no mesmo passo temporal global; a grade muda a cada \textit{beat}. Obstáculos são estáticos no tempo.

\subsection{Discreto versus contínuo}

O espaço de estados e o conjunto de ações são \textbf{discretos} (células e movimentos cardinais). Não há variáveis contínuas de controle.

\section{IMPLEMENTAÇÃO}

\subsection{Repositório e licença de uso}

O código-fonte completo encontra-se em repositório público, com histórico de versões e instruções de execução:

\begin{center}
\url{https://github.com/robsonpaulista/posicevagentes}
\end{center}

O material destina-se ao uso acadêmico no âmbito da disciplina.

\subsection{Requisitos e execução}

Exige-se \textbf{Python 3.10 ou superior}. Não há instalação obrigatória de pacotes via \texttt{pip}; utiliza-se a biblioteca padrão, incluindo \texttt{urllib} para chamadas HTTPS à API Groq e, opcionalmente, \texttt{tkinter} para interface gráfica (em distribuições Linux, pode ser necessário o pacote \texttt{python3-tk} do sistema).

Na raiz do repositório, recomenda-se:

\begin{verbatim}
python main.py
python main.py --visual --steps 120
\end{verbatim}

Para uso do LLM, copia-se \texttt{.env.example} para \texttt{.env} e define-se \texttt{GROQ\_API\_KEY}. Na ausência de chave ou em falha de rede ou cota, o sistema ativa \textbf{modo demonstração}, preservando a estrutura de decisão e o formato JSON de formação.

\subsection{Arquitetura dos agentes}

O \textbf{agente coreógrafo} implementa deliberação de alto nível com LLM: \textit{prompt} de sistema fixo exige resposta em JSON (\texttt{formation}, \texttt{center\_row}, \texttt{center\_col}, \texttt{scale}, \texttt{note}). O \textbf{agente piloto} mantém modelo interno da grade; para cada drone, calcula passo com BFS em grade, considerando obstáculos, posições alheias e intenções já planejadas de drones com índice menor, alinhando alvos por proximidade para reduzir bloqueios por permutação.

\section{CONCLUSÃO}

A simulação cumpre os requisitos do trabalho avaliativo: ambiente original implementado em Python, dois agentes autônomos com papéis distintos, integração de LLM com engenharia de \textit{prompt} e agente clássico baseado em modelo. A documentação no repositório e o presente relatório atendem às exigências de especificação PEAS, classificação dimensional e instruções de reprodutibilidade. Como extensões futuras, caberia modelar incerteza parcial ou custos de bateria, e avaliar formalmente a otimalidade do emparelhamento drone--alvo.

\newpage
\begin{thebibliography}{99}

\bibitem{russell2021}
RUSSELL, S.; NORVIG, P. \textbf{Inteligência artificial: uma abordagem moderna}. 4. ed. São Paulo: Pearson, 2021.

\bibitem{groq2024}
GROQ. \textbf{GroqCloud API documentation}. Disponível em: \url{https://console.groq.com/docs}. Acesso em: 5 abr. 2026.

\end{thebibliography}

\end{document}
